\documentclass[11pt]{article}
\usepackage[a4paper,margin=27mm]{geometry}
\usepackage{amsmath,amssymb,mathtools}
\usepackage{fontspec,unicode-math}
\setmainfont{Latin Modern Roman}
\setmathfont{Latin Modern Math}
\usepackage{xurl}
\usepackage[hidelinks]{hyperref}
\newcommand{\sref}[2]{\hyperlink{source:#1:#2}{[S#2]}}
\newcommand{\cataloguescope}[1]{\paragraph{Mathematical question.}#1}
\hypersetup{pdftitle={Problem statement}}
\begin{document}
\section{Pollock's tetrahedral-number conjecture}\label{statement}

\subsection{Definitions and mathematical statement}

\paragraph{Shared notation.}$\mathbb N=\{1,2,\ldots\}$, and $\mathbb Z,\mathbb Q,\mathbb R,\mathbb C$ denote the integers, rationals, reals and complex numbers. Any other conventions are specified locally.


For $n\in\mathbb Z_{\ge0}$ define the tetrahedral number $T_n=n(n+1)(n+2)/6$. The term $T_0=0$ permits a representation with fewer than five terms to be padded to five; repetitions are permitted.
\paragraph{Statement recorded in the supplied source.}
For every $N\in\mathbb N$, do there exist $n_1,\ldots,n_5\in\mathbb Z_{\ge0}$ such that $N=T_{n_1}+\cdots+T_{n_5}$? \sref{statement}{2}
\subsection{Short English statement}
Can every positive integer be written as a sum of at most five tetrahedral numbers?
\subsection{Sources}
\begin{description}
\item[\hypertarget{source:statement:1}{S1}] ULAM.AI and the UnsolvedMath Contributors, UnsolvedMath: A Curated Collection of Open Mathematics Problems, catalogue entry. CC BY 4.0. \url{https://huggingface.co/datasets/ulamai/UnsolvedMath}.
\item[\hypertarget{source:statement:2}{S2}] Zarathustra Brady, \emph{Sums of seven octahedral numbers}, arXiv:1509.04316v2 (2015) Introduction, p. 1, Pollock conjectures. \url{https://arxiv.org/pdf/1509.04316}.
\item[\hypertarget{source:statement:3}{S3}] Wikipedia contributors, \emph{List of unsolved problems in mathematics}, revision 1366636767 (2026) Number theory, Pollock's tetrahedral-number conjecture. \url{https://en.wikipedia.org/w/index.php?title=List_of_unsolved_problems_in_mathematics&oldid=1366636767#Number_theory}.
\end{description}
\label{end:statement}
\end{document}
